\documentclass[11pt]{article}
\usepackage[margin=25mm]{geometry}
\usepackage[T1]{fontenc}
\usepackage{booktabs}
\usepackage{array}
\usepackage{graphicx}
\usepackage{xurl}

\title{Processing Blue 2, 3, and 4}
\date{17 August 2026}

\begin{document}
\maketitle

\section{Capture data}

The underwater acoustic channel exhibits extended multipath and rapid time
variation.  We process the native Blue 2, Blue 3, and Blue 4 captures without
changing their measured channel records.

Each capture has three hydrophone signals.  The modem sampling rate is
4882.8125 samples/s, the center frequency is 13 kHz, and the snapshot rate is
48.828125 snapshots/s.  We use the first 47 s of each capture.

We verify the capture files against the secure hash algorithm (SHA-256) values
before processing.  The full values recorded in the result manifest are:

\begin{center}
\small
\begin{tabular}{@{}lll@{}}
\toprule
Capture & Bytes & SHA-256 \\
\midrule
\texttt{blue\_2.mat} & 206147418 &
\url{a3b7452c6c999777b49e5b4ddb3ef113ba4f40ff2d1504847f6d4be63e6a4bee} \\
\texttt{blue\_3.mat} & 206147418 &
\url{e126e95880513453722b8e9530bd55347fdb1605b76210f99375d1226b4acf9f} \\
\texttt{blue\_4.mat} & 206132594 &
\url{24eaa514da13aff5b5f3f7f8c81eb0fd0dd664ee5ee009b450169c63999c766b} \\
\bottomrule
\end{tabular}
\end{center}

\section{Replay}

The measured channel replay applies the measured phase and delay records to
the transmitted waveform.  The \texttt{theta\_hat} record applies phase,
whereas the \texttt{phi\_hat} record applies phase and delay.

We use guarded cubic interpolation for the delay record.  The replay retains
boundary samples with adaptive guards, requires all samples to remain inside
the measured channel and tracking support, and rejects non-monotone time maps.

We add independent proper-complex additive white Gaussian noise after the
measured replay and before receiver alignment.  Hence the signal-to-noise
power ratio does not alter the measured channel record.

\section{Acquisition}

The synchronization waveform occurs before and after the orthogonal frequency
division multiplexing (OFDM) body.  We require both occurrences and both chirp
components to provide credible synchronization evidence.

We use the opposite chirp slopes to separate carrier offset from duration
scale.  We reject a zero-energy waveform, missing peaks, implausible spacing,
an incomplete OFDM body, and carrier candidates that do not retain credible
synchronization.

The phase difference between the repeated synchronization waveforms defines a
bounded carrier-offset set between $-15$ Hz and 15 Hz.  We compare the
candidates using one common cyclic-prefix timing map, outer-pilot response, and
the cyclic-prefix residual.

We remove the selected carrier rotation over the complete waveform and then
repeat synchronization.  This order prevents a candidate-specific timing
change from deciding the carrier offset.

\section{Standard receiver}

Sparse quadrature phase shift keying (QPSK) pilots may not describe a phase
change between adjacent pilots.  However, the fourth power removes the unknown
QPSK data symbol and gives the channel phase modulo $\pi/2$ at every active
carrier.

We resolve the four possible phase states against the known outer pilots.  The
transition cost is weighted by received magnitude, so a phase transition costs
less at a channel null than at a strong carrier.

We apply this candidate only to frame-wide QPSK with at least two outer pilots
and a maximum pilot gap of at least 16 active carriers.  Forward error
correction (FEC) keeps the better of this candidate and the pilot-interpolated
candidate.

Although this change improves the Standard OFDM and FEC result, it is not used
by the partial-FFT or iterative receiver paths.  Their results must therefore
be checked separately and must not be inferred from the Standard result.

\section{Configuration}

The completed configuration uses $N=512$, a cyclic prefix of 64 samples, code
rate 0.25, outer-pilot spacing 6, inner-pilot spacing 8, check degree 14, and
four partial-FFT parts.  The horizon is filled by the available frame budget.

We use signal-to-noise power ratios from 0 dB to 30 dB in 2 dB steps, random
number generator seed 4, and 32 frames at each point.  Each frame contains 726
payload bits.

The frame contains 4388 samples and lasts 0.8986624 s.  This duration is within
the one-second frame budget at the native modem sampling rate.

We ran six receivers: \texttt{cwz\_joint}, \texttt{direct\_cz},
\texttt{lite}, \texttt{ofdm\_fec}, \texttt{pfft}, and
\texttt{profiled\_cz}.  The two cyclic redundancy check (CRC) paths return the
gradient result only when its CRC is valid.

We reran only Blue 2, Blue 3, and Blue 4, with three hydrophones for each
capture.  The three Blue 1 paths were retained and their files were checked
against the SHA-256 values recorded before the rerun.

\section{Validation}

The completed result records source commit
\texttt{5e518a4e62bb29df654f42b2b37d25c49d2c2524} and acquisition source SHA-256
\url{3679bf6f490dda5d307924405f0b404e8c7147b15603625c8daa252889c0e243}.

We require 12 path contracts, 1152 aggregate rows, 36864 frame-trace rows,
12288 selection-trace rows, 12 panels, and 72 receiver series.  Each refreshed
Blue 2--Blue 4 path contains 96 aggregate rows, 3072 frame rows, and 1024
selection rows.

We also require every refreshed path to have a Standard OFDM bit error rate
(BER) below 0.1 at 30 dB.  The maximum measured value over the nine paths is
0.0155819559, and all four worker error files are empty.

We include only a configuration with its completion marker and successful
independent validation.  Results from configurations still being processed are
not included in this note.

\section{Results}

Table~\ref{tab:ber30} gives the measured BER at 30 dB for the nine refreshed
paths.  A table entry of zero denotes no observed errors in 23232 payload bits;
it does not denote zero error probability.

\begin{table}[htbp]
\centering
\caption{Measured BER at 30 dB for the completed $N=512$ configuration.}
\label{tab:ber30}
\resizebox{\textwidth}{!}{%
\begin{tabular}{@{}llrrrrrr@{}}
\toprule
Capture & Hydrophone & \texttt{cwz\_joint} & \texttt{direct\_cz} &
\texttt{lite} & \texttt{ofdm\_fec} & \texttt{pfft} &
\texttt{profiled\_cz} \\
\midrule
Blue 2 & 1 & 0.000129 & 0.000129 & 0.000129 & 0 & 0.058196 & 0.000129 \\
Blue 2 & 2 & 0.008781 & 0.009900 & 0.008781 & 0 & 0.073218 & 0.008781 \\
Blue 2 & 3 & 0.000043 & 0.000043 & 0.000043 & 0 & 0.077565 & 0.000043 \\
Blue 3 & 1 & 0.380811 & 0.410253 & 0.380811 & 0 & 0.404787 & 0.380811 \\
Blue 3 & 2 & 0.389032 & 0.413137 & 0.389032 & 0 & 0.401946 & 0.389032 \\
Blue 3 & 3 & 0.380467 & 0.409392 & 0.380467 & 0 & 0.398244 & 0.380467 \\
Blue 4 & 1 & 0.184100 & 0.206999 & 0.184100 & 0 & 0.281939 & 0.184100 \\
Blue 4 & 2 & 0.156207 & 0.177514 & 0.156207 & 0.015582 & 0.269628 & 0.156207 \\
Blue 4 & 3 & 0.168991 & 0.187457 & 0.168991 & 0 & 0.300749 & 0.168991 \\
\bottomrule
\end{tabular}
}
\end{table}

The Standard OFDM result satisfies the stated 30 dB validation condition on
all nine refreshed paths.  However, the other receiver results remain poor on
Blue 3 and Blue 4 under this completed configuration.

Hence this evidence validates the Standard receiver result for the stated
condition.  It does not validate the partial-FFT, Direct C,z, JUNA-Lite, or
CRC-gated receiver results on Blue 3 and Blue 4.

\section{Sources}

We obtained the processing behavior from
\path{JunaCore/src/juna/common.jl},
\path{JunaCore/src/juna/frame_wide_ldpc.jl}, and
\path{JunaCore/tools/explorer/chain.json}.  We obtained the configuration and
counts from the completed result manifest and the 12 path contracts.

We checked completion with
\path{JunaCore/bench/validate_blue_native_awgn_direct_cz_extension.py} and
\path{JunaCore/bench/validate_blue_cfofix_channels234.py}.  The full measured
BER values are stored in the completed aggregate CSV file.

\end{document}
